\subsubsection{Subproceso 1: Acreditaci\'on y Pre-requisitos}

El modelo BPMN del Subproceso 1, con los actores definidos en la secci\'on 3, organiza las actividades del proceso actual en cinco carriles verticales. Los actores son: el Estudiante Regular, la Unidad de Tesorer\'ia / Caja UNI y Dependencias de Facultad, la Comisi\'on de Matr\'icula, Escuela Profesional y Tutores, la Direcci\'on de Registro Central y Estad\'istica (DIRCE / SIGA) y el Portal Intranet DIRCE / Intralu. La Caja UNI y las dependencias de Facultad registran los pagos del autoseguro y el estado de adeudos. La Escuela Profesional y los tutores conducen las gestiones acad\'emicas previas: el cierre de actas de evaluaci\'on, la consejer\'ia obligatoria a estudiantes en riesgo acad\'emico y el env\'io de n\'ominas de egresantes y convalidaciones. La DIRCE consolida la informaci\'on en el SIGA y valida los requisitos normativos. Tambi\'en ejecuta el c\'alculo de promedios y orden de m\'erito. Asigna los turnos y genera las fichas. Por su parte, la Intranet DIRCE publica los resultados y habilita el acceso en el horario establecido.

Esta vista vertical reproduce el flujo secuencial destacando las responsabilidades y compuertas de decisi\'on. Si el estudiante arrastra cursos desaprobados dos o m\'as veces, requiere obligatoriamente el pase de su tutor para activar su c\'odigo. Si tramita un cambio de plan o convalidaci\'on, la Facultad debe remitir el expediente f\'isico a la DIRCE para su registro manual en el sistema, lo que representa un cuello de botella recurrente. La validaci\'on determinante es el aporte del autoseguro m\'edico. Si no figura registrado, el acceso a la matr\'icula se bloquea. De estar conforme, el sistema asigna el turno correspondiente y genera la Ficha de Matr\'icula individual. El proceso finaliza cuando el portal despliega los cuatro controles de habilitaci\'on y deja al estudiante expedito para iniciar el Subproceso~2.

\begin{figure}[H]
    \centering
    \includegraphics[width=0.9\textwidth]{asis_subproceso_1_bpmn.png}
    \caption{Modelo BPMN vertical (AS IS) del Subproceso 1}
    \label{fig:asis_1_bpmn}
\end{figure}
